\section{Discussion}\label{sec:discussion}

The compactification constructed here has a local region in which
one can compare actual HE geometry with two physically identified
quantities. The leading vector hierarchy is controlled relatively
even when its smallest scale is arbitrarily weaker than the
others. The defined first-order tree pencil includes the leading
metric and two-form vector towers, and its ASD cancellation
preserves the full unbroken kernel. The scalar calculation uses
the same family, but requires derivative estimates and a compensated
reduction to identify a leading ambient interaction coefficient.
Its angular restriction is essential.

The two results have different uniformity. The spectral theorem
uses energy-relative reconstruction and survives changing orbit
scales without differentiating a horizontal projector. Curvature
uses smooth holomorphic regularity on a fixed angular patch.
The explicit projector derivatives and nonzero cone curvature
show why a uniform derivative assertion needs further information.
A quantitative lower bound on the admissible holomorphic chart
scale would be one route toward such a result.

The hyperk\"ahler quotient, qualitative Higgs mechanism,
heterotic bundle kinetic term, Hodge decomposition and
curvature--amplitude relation are established ingredients. The
substantial analytic premise is the local relative comparison
in \cite{wuebben15}; the present contribution is its normalized,
coupled use in an actual compactification with a worked interaction
and explicit error regimes. The local harmonic argument from
\cite{wuebben20} supplies the needed coordinates without a global
claim for the rank-six moduli space.

Several stronger questions remain. An exact finite-$\alpha'$
background with numerical remainder estimates would be needed
to turn the vector pencil bound into an exact corrected-pole
bound. A simultaneous interaction theorem would require control
of higher physical coefficients and their derivatives near
enhancement. Finite-coupling corrections require separate estimates.
Nonzero-FI wall crossing, bubbling and quantum amplitudes are
also outside the present result. At energies above the smallest
Higgs mass, the three massive multiplets must be restored rather
than extending the quotient sigma model beyond its window.
